% ═════════════════════════════════════════════════════════════════════════════
% CHAPTER 1: General Introduction
% ═════════════════════════════════════════════════════════════════════════════

\chapter{General Introduction}
\label{chap:introduction}

Increasing time pressure within the vehicle development process also leads to the ever greater importance of using engineering resources effectively. Furthermore, OEMs are increasingly working on numerous vehicle programs in parallel,  which requires managing multiple sophisticated disciplines ranging from embedded software development and control engineering to packaging design.  To keep pace with this increasing complexity and speed up vehicle development time and time to market,  automation of the decision support processes becomes more and more important due to the digitization in engineering.

\vspace{0.3cm}

This internship was carried out within the Mechatronics Engineering department within PSA's Research \& Development division, located in Sochaux, East of France. I was the second engineer I know who came to this department,  and the only one in the Mechatronics division.  This department was dedicated to the integration of the electrical and electromechanical parts of all vehicles assembled in PSA's plant of Sochaux.  The department dealt with projects in all their different statuses,  appropriation means and priority levels: some projects were fully finished and delivered, while some others were still in their development phase.  The main difficulty of the Mechatronics division was the allocation of tasks and resources. No formal assignment and distribution system was in place,  hence the project managers' reliance on their experience.

\vspace{0.3cm}

The limitations of the current approach can be summarized as follows; Firstly, the workload and capacity information will never be perfectly accurate. Secondly,  and related, it takes considerable,  effort to locate the engineer best suited for a specific task.  Thirdly,  the current approach is unable to identify a task that cannot be allocated based upon interdependencies or the lack of required skills in the pool of available staff, until it is at the allocation phase. Fourthly, the slow manual validation of design specifications and the lack of systematic checks against Stellantis engineering standards leads to rework cycles and potential compliance issues.

\vspace{0.3cm}

This system support the governance of Mechatronics in several ways. First, it aggregates the team's information about individuals and resources. Second,it supportshuman-based decision making by automatically presenting information to administrators and project managers at the time their assistance is needed. It also supports the automatic validation and checking of assignments in areas such as skill sets or availability by checking individual characteristics and project requirements against rules on assignments.The specific goals for this work are to create an AI assistant agent for the different daily needs of mecatronics team.

\vspace{0.3cm}
This report follows a structured approach to address the research question through successive analytical phases:

\begin{itemize}
  \item \textbf{Chapter 1 -- State of the Art:} Reviews existing resource management approaches, intelligent agent systems, constraint satisfaction techniques, and knowledge representation methods, establishing the theoretical foundation.
  
  \item \textbf{Chapter 2 -- Current System Analysis:} Describes the Mecatronics team's current resource allocation process, existing tools, constraints, and gaps requiring automation.
  
  \item \textbf{Chapter 3 -- Requirements Specification:} Presents formal functional and non-functional requirements for the proposed system and its components.
  
  \item \textbf{Chapter 4 -- System Design:} Proposes the intelligent agent architecture with components for competency discovery, workload analysis, constraint satisfaction, specification validation, and access control.
  
  \item \textbf{Chapter 5 -- Implementation and Prototyping:} Describes development process, tools, integration approach, and the resulting prototype system.
  
  \item \textbf{Chapter 6 -- Evaluation and Results:} Presents testing protocols, performance metrics, validation against baseline methods, and lessons learned with future research directions.
\end{itemize}

The methodology combines requirements analysis, knowledge engineering, constraint programming, and iterative prototyping, with validation conducted against real project scenarios from the Mechatronics team.
